
反德西特空间（Anti-de Sitter space, AdS）是具有常负曲率的时空，它是 AdS/CFT 对偶中体空间的几何背景。AdS 空间在理论物理学中有着深远的历史：它最早由物理学家在研究宇宙学常数为负的爱因斯坦场方程时发现，后来在弦理论和全息对偶的研究中扮演了核心角色。理解 AdS 空间的几何性质对于理解全息对偶至关重要。

本节将从 AdS 空间的嵌入定义出发，系统地介绍其各种坐标系、几何性质、因果结构、测地线、波动方程以及黑洞热力学。我们将看到，AdS 空间丰富的几何结构为 AdS/CFT 对偶提供了坚实的基础。

%=============================================================================
\section{de Sitter 空间与 Anti-de Sitter 空间的对比}

在介绍 AdS 空间之前，有必要先理解它与 de Sitter 空间（dS）的关系。这两者都是具有最大对称性的爱因斯坦方程解，区别在于宇宙学常数的符号。

\begin{definition}[最大对称时空]
一个 $d+1$ 维时空称为\textbf{最大对称}（maximally symmetric）的，如果它具有最大可能数目的等距变换，即 $(d+1)(d+2)/2$ 个。最大对称时空的黎曼张量完全由里奇标量决定：
\begin{equation}
    R_{\mu\nu\rho\sigma} = \frac{R}{d(d+1)} \left( g_{\mu\rho} g_{\nu\sigma} - g_{\mu\sigma} g_{\nu\rho} \right).
\end{equation}
对于爱因斯坦方程 $R_{\mu\nu} - \frac{1}{2} g_{\mu\nu} R + \Lambda g_{\mu\nu} = 0$，最大对称解满足 $R_{\mu\nu} = \frac{2\Lambda}{d-1} g_{\mu\nu}$。
\end{definition}

\begin{table}[htbp]
\centering
\begin{tabular}{|l|c|c|c|}
\hline
\textbf{时空} & \textbf{宇宙学常数 $\Lambda$} & \textbf{里奇标量 $R$} & \textbf{等距群} \\
\hline
闵可夫斯基 & $0$ & $0$ & $ISO(1,d)$ \\
de Sitter (dS) & $> 0$ & $> 0$ & $SO(1,d+1)$ \\
Anti-de Sitter (AdS) & $< 0$ & $< 0$ & $SO(2,d)$ \\
\hline
\end{tabular}
\caption{三种最大对称时空的对比。AdS 空间具有负宇宙学常数和负里奇标量。}
\label{tab:maximally_symmetric}
\end{table}

de Sitter 空间描述了一个具有正宇宙学常数的膨胀宇宙，它是当前宇宙加速膨胀的近似模型。Anti-de Sitter 空间则描述了一个具有负宇宙学常数的"收缩"宇宙。虽然 AdS 空间不是我们宇宙的好的近似（我们的宇宙具有正宇宙学常数），但它在 AdS/CFT 对偶中扮演了核心角色，因为它的边界结构特别适合定义共形场论 \cite{Maldacena1997}。

\begin{remark}[AdS 空间在弦理论中的起源 \cite{Maldacena1997}]
AdS 空间在现代理论物理中的重要性源于 Maldacena 的发现：IIB 型弦理论在 AdS$_5 \times S^5$ 背景上等价于定义在 AdS 边界上的 $\mathcal{N}=4$ 超对称 Yang-Mills 理论 \cite{Maldacena1997}。具体来说，考虑 $N$ 个 D3 膜的低能有效理论：
\begin{itemize}
    \item \textbf{开弦描述}（低能极限）：$U(N)$ $\mathcal{N}=4$ SYM 理论
    \item \textbf{闭弦描述}（低能极限）：AdS$_5 \times S^5$ 上的 IIB 型弦理论
\end{itemize}
当 $g_s N \gg 1$（即 $\lambda = g_{\rm YM}^2 N \gg 1$）时，弦理论的 $\alpha'$ 修正可以忽略，AdS 的曲率半径 $L$ 远大于弦长 $l_s$，此时弦理论经典地描述了强耦合的 SYM 理论。
\end{remark}

%=============================================================================
\section{AdS 空间的定义与嵌入构造}

AdS 空间可以通过嵌入在更高维平坦空间中的超曲面来定义。这种定义方式不仅清晰地显示了 AdS 的等距群，而且为推导各种坐标系提供了便利 \cite{Aharony1999}。

\begin{definition}[AdS 空间的嵌入定义]
AdS$_{d+1}$ 空间可以定义为嵌入在具有号差 $(2,d)$ 的平坦空间 $\mathbb{R}^{2,d}$ 中的超曲面。嵌入空间的度规为：
\begin{equation}
    ds^2_{\rm ambient} = -dX_{-1}^2 - dX_0^2 + dX_1^2 + \cdots + dX_d^2,
\end{equation}
AdS$_{d+1}$ 是满足以下约束的超曲面：
\begin{equation}
    -X_{-1}^2 - X_0^2 + X_1^2 + \cdots + X_d^2 = -L^2,
    \label{eq:AdS_embedding}
\end{equation}
其中 $L > 0$ 是 AdS 半径（也称为曲率半径），它决定了 AdS 空间的"大小"和曲率。
\end{definition}

\begin{remark}
约束方程 \eqref{eq:AdS_embedding} 可以写成更紧凑的形式。定义嵌入坐标 $X^A$（$A = -1, 0, 1, \ldots, d$），以及 $\mathbb{R}^{2,d}$ 上的度规 $\eta_{AB} = \text{diag}(-1,-1,+1,\ldots,+1)$，则约束为：
\begin{equation}
    \eta_{AB} X^A X^B = X^A X_A = -L^2.
\end{equation}
这里我们使用爱因斯坦求和约定，$X_A = \eta_{AB} X^B$。
\end{remark}

\subsection{等距群的推导}

AdS$_{d+1}$ 的等距群可以从嵌入定义直接读出。约束方程 \eqref{eq:AdS_embedding} 在 $\mathbb{R}^{2,d}$ 的线性变换 $X^A \to {M^A}_B X^B$ 下保持不变，当且仅当 $M$ 满足：
\begin{equation}
    M^T \eta M = \eta, \quad \eta = \text{diag}(-1,-1,+1,\ldots,+1).
\end{equation}
这正是群 $O(2,d)$ 的定义。由于我们要求变换保持时间方向（即不翻转 $X_{-1}$ 和 $X_0$ 的符号），等距群为 $SO(2,d)$（或更精确地说是其单位连通分支 $SO^+(2,d)$）。

$SO(2,d)$ 的维数为：
\begin{equation}
    \dim SO(2,d) = \frac{(d+2)(d+1)}{2} - 1 = \frac{d(d+3)}{2}.
\end{equation}
对于 $d=4$（即 AdS$_5$），等距群有 $\frac{4 \times 7}{2} = 14$ 个生成元。这恰好等于 4 维共形群 $SO(2,4)$ 的维数，这是 AdS/CFT 对偶的对称性基础 \cite{Maldacena1997}。

\begin{remark}[等距群与共形群的对应 \cite{Aharony1999}]
AdS$_{d+1}$ 的等距群 $SO(2,d)$ 恰好是 $d$ 维闵可夫斯基时空的共形群。这不是巧合——它是 AdS/CFT 对偶的核心对称性结构。等距群的生成元 $J_{AB}$ 可以一一映射到共形群的生成元：
\begin{equation}
    J_{ij} \leftrightarrow M_{ij}, \quad J_{-1,\mu} \leftrightarrow P_\mu, \quad J_{0,\mu} \leftrightarrow K_\mu, \quad J_{-1,0} \leftrightarrow D,
\end{equation}
其中 $M_{ij}$ 是旋转，$P_\mu$ 是平移，$K_\mu$ 是特殊共形变换，$D$ 是标度变换 \cite{Aharony1999}。这意味着 AdS 空间中的对称性操作在边界上表现为共形变换，这是全息对偶的对称性基础。
\end{remark}

\subsection{等距群的生成元}

$SO(2,d)$ 的生成元 $J_{AB}$（$A, B = -1, 0, 1, \ldots, d$）满足李代数：
\begin{equation}
    [J_{AB}, J_{CD}] = i(\eta_{AD} J_{BC} - \eta_{AC} J_{BD} + \eta_{BC} J_{AD} - \eta_{BD} J_{AC}).
\end{equation}
在嵌入空间中，生成元可以表示为：
\begin{equation}
    J_{AB} = i(X_A \partial_B - X_B \partial_A).
\end{equation}
这些生成元可以分为以下几类：
\begin{itemize}
    \item \textbf{旋转} $J_{ij}$（$i, j = 1, \ldots, d$）：对应于边界空间中的旋转，生成元为 $M_{\mu\nu} = i(x_\mu \partial_\nu - x_\nu \partial_\mu)$
    \item \textbf{平移} $J_{-1,i}$ 和 $J_{0,i}$：对应于边界空间中的平移，生成元为 $P_\mu = -i\partial_\mu$
    \item \textbf{标度变换} $J_{-1,0}$：对应于边界上的标度变换，生成元为 $D = -ix^\mu \partial_\mu$
    \item \textbf{特殊共形变换}：对应于边界上的特殊共形变换，生成元为 $K_\mu = i(x^2 \partial_\mu - 2x_\mu x^\nu \partial_\nu)$
\end{itemize}

%=============================================================================
\section{AdS 的坐标系}

AdS 空间有多种坐标系，每种坐标系都有其特殊的用途。我们将详细介绍几种最重要的坐标系，并展示它们之间的变换关系 \cite{Aharony1999}。

\subsection{全局坐标（Global Coordinates）}

全局坐标覆盖整个 AdS 空间，是最自然的坐标系。

\begin{definition}[全局坐标]
全局坐标 $(t, r, \Omega)$ 定义为：
\begin{align}
    X_{-1} &= L \frac{\cos t}{\cos \rho}, \notag \\
    X_0 &= L \frac{\sin t}{\cos \rho}, \notag \\
    X_i &= L \tan \rho \, \hat{n}_i, \quad i = 1, \ldots, d,
\end{align}
其中 $\hat{n}_i$ 是 $S^{d-1}$ 上的单位矢量（$\sum_i \hat{n}_i^2 = 1$），$\rho \in [0, \pi/2)$，$t \in [0, 2\pi)$。
\end{definition}

在全局坐标下，AdS 度规为：
\begin{equation}
    ds^2 = \frac{L^2}{\cos^2 \rho} \left( -dt^2 + d\rho^2 + \sin^2 \rho \, d\Omega_{d-1}^2 \right).
    \label{eq:AdS_global_rho}
\end{equation}

定义 $r = L \tan \rho$（即 $r \in [0, \infty)$），度规变为：
\begin{equation}
    ds^2 = -\left(1 + \frac{r^2}{L^2}\right) dt^2 + \frac{dr^2}{1 + r^2/L^2} + r^2 d\Omega_{d-1}^2.
    \label{eq:AdS_global}
\end{equation}

\begin{remark}
在全局坐标下，AdS 空间的边界位于 $r \to \infty$（即 $\rho \to \pi/2$），边界拓扑为 $\mathbb{R} \times S^{d-1}$。时间坐标 $t$ 的周期为 $2\pi$（对于紧致化的 AdS）或 $-\infty < t < \infty$（对于非紧致化的 AdS）。在物理应用中，我们通常使用非紧致化的 AdS。
\end{remark}

\subsection{庞加莱坐标的推导}

庞加莱坐标是 AdS/CFT 对偶中最常用的坐标系。我们从嵌入定义出发，详细推导庞加莱坐标。

\textbf{物理动机}：在 AdS/CFT 对偶中，我们希望边界理论定义在闵可夫斯基时空中。这意味着我们需要找到一个坐标系，使得边界是平坦的闵可夫斯基时空，而不是弯曲的 $\mathbb{R} \times S^{d-1}$。庞加莱坐标正是这样的坐标系。

\textbf{推导思路}：我们从嵌入约束方程出发：
\begin{equation}
    -X_{-1}^2 - X_0^2 + X_1^2 + \cdots + X_d^2 = -L^2.
\end{equation}
我们希望找到一组坐标 $(t, z, \vec{x})$，使得：
\begin{enumerate}
    \item 边界（$z \to 0$）是闵可夫斯基时空 $\mathbb{R}^{d-1,1}$
    \item 度规具有共形平坦的形式 $ds^2 = \Omega^2(z)(-dt^2 + d\vec{x}^2 + dz^2)$
\end{enumerate}

\textbf{具体推导}：考虑以下坐标变换：
\begin{align}
    X_0 &= \frac{L t}{z}, \label{eq:Poincare_X0} \\
    X_i &= \frac{L x_i}{z}, \quad i = 1, \ldots, d-1, \label{eq:Poincare_Xi} \\
    X_{-1} + X_d &= \frac{L^2}{z}. \label{eq:Poincare_sum}
\end{align}
这些定义的物理意义是：
\begin{itemize}
    \item $X_0$ 和 $X_i$ 与边界坐标 $t$ 和 $x_i$ 成正比，比例因子为 $L/z$
    \item $X_{-1} + X_d$ 与径向坐标 $z$ 成反比
\end{itemize}

从 \eqref{eq:Poincare_sum} 可以解出 $X_d$：
\begin{equation}
    X_d = \frac{L^2}{z} - X_{-1}.
\end{equation}
代入嵌入约束方程：
\begin{equation}
    -X_{-1}^2 - \frac{L^2 t^2}{z^2} + \frac{L^2 \vec{x}^2}{z^2} + \left(\frac{L^2}{z} - X_{-1}\right)^2 = -L^2.
\end{equation}
展开并整理：
\begin{equation}
    -X_{-1}^2 - \frac{L^2 t^2}{z^2} + \frac{L^2 \vec{x}^2}{z^2} + \frac{L^4}{z^2} - \frac{2L^2 X_{-1}}{z} + X_{-1}^2 = -L^2.
\end{equation}
$X_{-1}^2$ 项抵消，得到：
\begin{equation}
    \frac{L^2}{z^2}\left(-t^2 + \vec{x}^2 + L^2\right) - \frac{2L^2 X_{-1}}{z} = -L^2.
\end{equation}
解出 $X_{-1}$：
\begin{equation}
    X_{-1} = \frac{z}{2}\left(1 + \frac{1}{z^2}(L^2 + \vec{x}^2 - t^2)\right).
\end{equation}
类似地，可以求出 $X_d$：
\begin{equation}
    X_d = \frac{z}{2}\left(1 - \frac{1}{z^2}(L^2 - \vec{x}^2 + t^2)\right).
\end{equation}

\textbf{验证约束}：将这些表达式代入嵌入约束方程，可以验证：
\begin{equation}
    -X_{-1}^2 - X_0^2 + X_1^2 + \cdots + X_d^2 = -L^2.
\end{equation}
这表明我们的坐标变换是正确的。

\textbf{计算度规}：现在我们来计算庞加莱坐标下的度规。嵌入空间的度规为：
\begin{equation}
    ds^2 = -dX_{-1}^2 - dX_0^2 + dX_1^2 + \cdots + dX_d^2.
\end{equation}
计算各微分：
\begin{align}
    dX_0 &= \frac{L}{z} dt - \frac{L t}{z^2} dz, \\
    dX_i &= \frac{L}{z} dx_i - \frac{L x_i}{z^2} dz, \\
    dX_{-1} &= \frac{1}{2}\left(1 - \frac{1}{z^2}(L^2 + \vec{x}^2 - t^2)\right) dz + \frac{1}{z}(x_i dx_i - t dt), \\
    dX_d &= \frac{1}{2}\left(1 + \frac{1}{z^2}(L^2 - \vec{x}^2 + t^2)\right) dz - \frac{1}{z}(x_i dx_i - t dt).
\end{align}
代入度规表达式，经过详细计算（这里省略中间步骤），得到：
\begin{equation}
    ds^2 = \frac{L^2}{z^2}\left(-dt^2 + d\vec{x}^2 + dz^2\right).
\end{equation}
这就是庞加莱坐标下的 AdS 度规！

\begin{remark}
庞加莱坐标的推导展示了如何从嵌入定义出发，通过巧妙的坐标变换得到所需的度规形式。关键的物理洞察是：为了使边界是闵可夫斯基时空，我们需要将 $X_0$ 和 $X_i$ 与边界坐标 $t$ 和 $x_i$ 成正比，而将 $X_{-1} + X_d$ 与径向坐标 $z$ 成反比。
\end{remark}

\begin{remark}
庞加莱坐标只覆盖 AdS 空间的一半（称为\textbf{庞加莱补丁}，Poincaré patch），因为 $z > 0$。边界位于 $z \to 0$ 处，边界拓扑为 $\mathbb{R}^{d-1,1}$（闵可夫斯基时空）。在 $z \to \infty$ 处，度规有一个视界（称为庞加莱视界），它将庞加莱补丁与 AdS 的其他部分隔开。
\end{remark}

\begin{remark}[庞加莱坐标与全息对偶 \cite{Witten1998}]
在 AdS/CFT 对偶的标准表述中 \cite{Witten1998}，庞加莱坐标是最常用的坐标系，原因如下：
\begin{itemize}
    \item 边界是平坦的闵可夫斯基时空，这对应于 CFT 定义在平坦时空上的标准设置
    \item 度规的共形平坦形式 $ds^2 = (L^2/z^2)(-dt^2 + d\vec{x}^2 + dz^2)$ 使得边界上的共形对称性显式可见
    \item 径向坐标 $z$ 自然地对应于 CFT 中的能量尺度：$z \sim 1/\mu$，边界 $z \to 0$ 对应于 UV，体空间内部 $z \to \infty$ 对应于 IR
\end{itemize}
然而，如果需要研究全局 AdS/CFT（即 CFT 定义在 $\mathbb{R} \times S^{d-1}$ 上），则必须使用全局坐标 \cite{Witten1998}。
\end{remark}

\subsection{从全局坐标到庞加莱坐标的变换}

全局坐标和庞加莱坐标之间的变换可以通过嵌入坐标来建立。具体来说，考虑以下变换：
\begin{align}
    t &= \frac{t_P}{\sqrt{1 - z_P^2/(4L^2)(1 + \vec{x}_P^2/L^2 - t_P^2/L^2)}}, \notag \\
    r &= \frac{\sqrt{1 + \vec{x}_P^2/L^2 - t_P^2/L^2}}{z_P/(2L)}.
\end{align}
这个变换将全局坐标 $(t, r)$ 映射到庞加莱坐标 $(t_P, z_P, \vec{x}_P)$。在边界附近（$r \to \infty$ 或 $z_P \to 0$），两种坐标系给出相同的边界结构。

\subsection{Fefferman-Graham 坐标}

Fefferman-Graham（FG）坐标是全息重整化中使用的坐标系，它使得边界附近的渐近展开更加清晰。

\begin{definition}[Fefferman-Graham 坐标]
对于渐近 AdS 时空，Fefferman-Graham 坐标 $\rho$ 定义为：
\begin{equation}
    ds^2 = \frac{L^2}{\rho^2} \left( d\rho^2 + g_{\mu\nu}(x, \rho) dx^\mu dx^\nu \right),
\end{equation}
其中 $g_{\mu\nu}(x, \rho)$ 在边界 $\rho \to 0$ 附近可以展开为：
\begin{equation}
    g_{\mu\nu}(x, \rho) = g_{(0)\mu\nu}(x) + \rho^2 g_{(2)\mu\nu}(x) + \cdots + \rho^d g_{(d)\mu\nu}(x) + h_{(d)\mu\nu}(x) \rho^d \ln \rho + \cdots
\end{equation}
\end{definition}

\begin{remark}
在纯 AdS 时空（庞加莱坐标）中，$g_{\mu\nu}(x, \rho) = \eta_{\mu\nu}$，即所有高阶修正为零。在一般的渐近 AdS 时空中，系数 $g_{(0)\mu\nu}$ 是边界 CFT 的背景度规，$g_{(d)\mu\nu}$ 与应力张量的期望值相关。这种展开是全息重整化的基础 \cite{deHaro2000, Balasubramanian1999}。
\end{remark}

\begin{remark}[共形边界 \cite{Aharony1999}]
AdS 空间的共形边界（conformal boundary）可以通过共形紧致化来定义 \cite{Aharony1999}。考虑 AdS$_{d+1}$ 的庞加莱度规：
\begin{equation}
    ds^2 = \frac{L^2}{z^2}\left(-dt^2 + d\vec{x}^2 + dz^2\right).
\end{equation}
做共形变换 $\tilde{g}_{\mu\nu} = (z/L)^2 g_{\mu\nu}$，得到：
\begin{equation}
    d\tilde{s}^2 = -dt^2 + d\vec{x}^2 + dz^2,
\end{equation}
这是 $(d+1)$ 维闵可夫斯基度规！在共形变换后的度规中，AdS 的边界 $z \to 0$ 变成了一个有限的超曲面。这个边界就是 CFT 所定义的时空。注意，共形变换改变了边界处的因果结构（从类时变为类光），但保持了共形结构不变。
\end{remark}

\subsection{Rindler 型坐标}

Rindler 型坐标对于理解 AdS 空间的因果结构和黑洞物理很重要。

\begin{definition}[Rindler 型坐标]
在 AdS 空间中，可以定义类似于闵可夫斯基时空中 Rindler 坐标的坐标系：
\begin{equation}
    ds^2 = -\frac{r^2}{L^2} dt^2 + \frac{L^2}{r^2} dr^2 + r^2 d\vec{x}^2.
\end{equation}
这个度规可以通过坐标变换 $z = L^2/r$ 从庞加莱坐标得到。
\end{definition}

%=============================================================================
\section{AdS 空间的几何性质}

AdS 空间具有丰富的几何结构，这些结构为 AdS/CFT 对偶提供了坚实的基础。我们来系统地介绍 AdS 的几何性质 \cite{Aharony1999}。

\subsection{曲率张量}

AdS$_{d+1}$ 空间是常负曲率空间。这个性质可以从嵌入定义直接推导出来。

\begin{property}[AdS 的曲率张量]
AdS$_{d+1}$ 空间的黎曼张量为：
\begin{equation}
    R_{\mu\nu\rho\sigma} = -\frac{1}{L^2} \left( g_{\mu\rho} g_{\nu\sigma} - g_{\mu\sigma} g_{\nu\rho} \right).
\end{equation}
从黎曼张量可以计算里奇张量和里奇标量：
\begin{align}
    R_{\mu\nu} &= -\frac{d}{L^2} g_{\mu\nu}, \\
    R &= -\frac{d(d+1)}{L^2}.
\end{align}
\end{property}

\begin{note}[曲率的物理意义]
里奇标量 $R = -d(d+1)/L^2$ 是负的，这反映了 AdS 空间的负曲率性质。
\begin{itemize}
    \item 对于 AdS$_5$（$d=4$），$R = -20/L^2$
    \item 曲率半径 $L$ 越大，曲率越小（即空间越"平坦"）
    \item 在 AdS/CFT 对偶中，$L$ 与边界 CFT 的中心荷相关：$c \sim L^{d-1}/G_N$
\end{itemize}
\end{note}

\subsection{截面曲率}

AdS 空间的\textbf{截面曲率}（sectional curvature）在所有二维切平面上都是常数 $-1/L^2$。这意味着 AdS 空间是\textbf{空间各向同性}的：在每一点，所有方向上的弯曲程度都相同。

\begin{definition}[截面曲率]
对于切空间中的两个线性无关的矢量 $u$ 和 $v$，截面曲率为：
\begin{equation}
    K(u, v) = \frac{R_{\mu\nu\rho\sigma} u^\mu v^\nu u^\rho v^\sigma}{(g_{\mu\rho} u^\mu u^\rho)(g_{\nu\sigma} v^\nu v^\sigma) - (g_{\mu\sigma} u^\mu v^\sigma)^2}.
\end{equation}
对于 AdS 空间，$K(u, v) = -1/L^2$ 对所有 $u, v$ 成立。
\end{definition}

\subsection{体积元与测地线距离}

AdS 空间的体积元为：
\begin{equation}
    \sqrt{-g} = \left(\frac{L}{z}\right)^{d+1}
\end{equation}
（在庞加莱坐标下）。这个因子在 $z \to 0$（边界附近）发散，反映了 AdS 空间在边界附近的无穷大体积。

两点之间的测地线距离（geodesic distance）$\sigma(x, y)$ 是 AdS 空间中重要的几何量。在嵌入空间中，两点 $X$ 和 $Y$ 之间的测地线距离为：
\begin{equation}
    \cosh\left(\frac{\sigma}{L}\right) = -\frac{X \cdot Y}{L^2} = \frac{X_{-1} Y_{-1} + X_0 Y_0 - X_1 Y_1 - \cdots - X_d Y_d}{L^2}.
\end{equation}

\begin{example}[庞加莱坐标中的测地线距离]
在庞加莱坐标中，两点 $(z, \vec{x})$ 和 $(z', \vec{x}')$ 之间的测地线距离为：
\begin{equation}
    \sigma = L \ln\left(\frac{(z - z')^2 + (\vec{x} - \vec{x}')^2}{2zz'} + \sqrt{\left(\frac{(z - z')^2 + (\vec{x} - \vec{x}')^2}{2zz'}\right)^2 - 1}\right).
\end{equation}
当两点都在边界附近（$z, z' \to 0$）时，测地线距离发散，这反映了 AdS 边界上的 UV 发散。
\end{example}

%=============================================================================
\section{AdS 的因果结构}

理解 AdS 空间的因果结构对于理解 AdS/CFT 对偶至关重要。AdS 空间有一个独特的性质：它的\textbf{类光无穷远是类时的}，这意味着光信号可以在有限时间内从边界传播到体空间内部并返回 \cite{Witten1998}。

\subsection{Penrose 图}

AdS$_{d+1}$ 的 Penrose 图（Penrose diagram）是一个矩形，如图 \ref{fig:AdS_Penrose} 所示。

\begin{figure}[htbp]
\centering
\begin{tikzpicture}[scale=1.5]
    % 绘制边界
    \draw[thick] (-1.5,-1.5) -- (-1.5,1.5) node[left] {$\mathscr{I}^+$};
    \draw[thick] (1.5,-1.5) -- (1.5,1.5) node[right] {$\mathscr{I}^-$};
    \draw[thick] (-1.5,1.5) -- (1.5,1.5) node[above] {$i^+$};
    \draw[thick] (-1.5,-1.5) -- (1.5,-1.5) node[below] {$i^-$};
    
    % 绘制 AdS 内部
    \fill[blue!10] (-1.5,-1.5) rectangle (1.5,1.5);
    
    % 绘制类光测地线
    \draw[red, thick, ->] (-1.2,-1.2) -- (0,0);
    \draw[red, thick, ->] (0,0) -- (1.2,1.2);
    
    % 标注
    \node at (0,0) {$i^0$};
    \node at (0,-1.8) {AdS 体空间};
    
    % 绘制时间方向
    \draw[->, thick] (0,-1.3) -- (0,1.3) node[midway, right] {$t$};
\end{tikzpicture}
\caption{AdS 空间的 Penrose 图。矩形区域代表整个 AdS 空间。左右边界 $\mathscr{I}^\pm$ 是类时的，上下顶点 $i^\pm$ 分别是未来和过去无穷远。}
\label{fig:AdS_Penrose}
\end{figure}

\begin{remark}
AdS 空间的 Penrose 图与闵可夫斯基时空的 Penrose 图有本质区别。在闵可夫斯基时空中，类光无穷远 $\mathscr{I}^\pm$ 是类光的（45度倾斜）；而在 AdS 空间中，类光无穷远是类时的（垂直）。这意味着：
\begin{itemize}
    \item 光信号可以在有限时间内从边界传播到体空间内部
    \item AdS 空间没有真正的"无穷远"——边界是有限距离的
    \item 这为 AdS/CFT 对偶提供了自然的边界条件
\end{itemize}
\end{remark}

\begin{remark}[AdS 的边界条件 \cite{Witten1998}]
Witten 在 \cite{Witten1998} 中详细讨论了 AdS 空间的边界条件问题。由于 AdS 的边界是类时的（而非类光的），边界上的场需要满足特定的边界条件才能唯一确定体空间中的解。对于质量为 $m$ 的标量场，边界上的渐近行为为：
\begin{equation}
    \phi(z, x) \sim z^{d-\Delta} \phi_0(x) + z^\Delta \langle \mathcal{O}(x) 
angle, \quad z \to 0,
\end{equation}
其中 $\Delta = \frac{d}{2} + \sqrt{\frac{d^2}{4} + m^2 L^2}$。Witten 指出，为了定义一个良定的变分问题，我们需要指定 $\phi_0(x)$（即 non-normalizable 模的边界值）作为边界条件 \cite{Witten1998}。这对应于在 CFT 中添加源项 $\int \phi_0 \mathcal{O}$。
\end{remark}

\begin{proof}[Penrose 图的构造]
我们来详细说明如何从全局坐标构造 Penrose 图。从全局度规出发：
\begin{equation}
    ds^2 = \frac{L^2}{\cos^2\rho}\left(-dt^2 + d\rho^2 + \sin^2\rho \, d\Omega_{d-1}^2\right),
\end{equation}
其中 $\rho \in [0, \pi/2)$，$t \in (-\infty, \infty)$。

引入新的坐标以将无穷远压缩到有限范围：
\begin{align}
    t + \rho &= \tan\left(\frac{T+R}{2}\right), \notag \\
    t - \rho &= \tan\left(\frac{T-R}{2}\right),
\end{align}
其中 $T \in (-\pi, \pi)$，$R \in [0, \pi)$。在 $(T, R)$ 坐标下，AdS 的因果结构变得完全透明：
\begin{itemize}
    \item 未来无穷远 $i^+$：$T = \pi$，$R = 0$
    \item 过去无穷远 $i^-$：$T = -\pi$，$R = 0$
    \item 空间无穷远 $i^0$：$R = \pi$，$T = 0$
    \item 类光无穷远 $\mathscr{I}^\pm$：$T = \pm(\pi - R)$
\end{itemize}
因此 Penrose 图是一个矩形，边界 $\mathscr{I}^\pm$ 是垂直的直线（类时的），这与闵可夫斯基时空（边界是倾斜的 45 度线）形成鲜明对比。
\end{proof}

\subsection{全局时间与因果性}

在全局坐标下，AdS 空间的时间坐标 $t$ 的范围是 $-\infty < t < \infty$（非紧致化版本）。由于类光无穷远是类时的，光信号可以在有限时间内从边界传播到体空间内部，然后返回边界。

考虑一个从边界点 $(t_0, r = \infty)$ 发出的类光信号。在全局坐标下，这个信号沿测地线传播，在时间 $\Delta t = \pi$ 后到达"对侧"边界。这意味着 AdS 空间的"往返时间"是有限的，为 $\pi L$（以全局时间计量）。

\subsection{庞加莱视界}

庞加莱坐标只覆盖 AdS 空间的一半。在 $z \to \infty$ 处，存在一个\textbf{庞加莱视界}（Poincaré horizon），它将庞加莱补丁与 AdS 的其他部分隔开。

\begin{remark}
庞加莱视界不是一个真正的因果视界——观察者可以穿过它。但它标志着庞加莱坐标的适用范围的边界。在 AdS/CFT 对偶中，庞加莱补丁对应于边界 CFT 定义在闵可夫斯基时空上的情况。如果需要覆盖整个 AdS 空间（例如研究全局 AdS/CFT），则需要使用全局坐标。
\end{remark}

%=============================================================================
\section{AdS 中的测地线}

测地线是弯曲时空中"直线"的推广，它描述了自由粒子在时空中运动的轨迹。在 AdS 空间中，测地线的行为对于理解全息对偶、黑洞物理和纠缠熵都很重要 \cite{Maldacena1997}。

\subsection{测地线方程}

在 AdS 空间中，测地线方程为：
\begin{equation}
    \frac{d^2 x^\mu}{d\tau^2} + \Gamma^\mu_{\rho\sigma} \frac{dx^\rho}{d\tau} \frac{dx^\sigma}{d\tau} = 0,
\end{equation}
其中 $\tau$ 是仿射参数，$\Gamma^\mu_{\rho\sigma}$ 是 Christoffel 符号。

\subsection{Christoffel 符号的推导}

在庞加莱坐标下，AdS 度规为 $ds^2 = (L^2/z^2)(-dt^2 + d\vec{x}^2 + dz^2)$。我们来详细推导 Christoffel 符号。

首先，写出度规张量的分量：
\begin{equation}
    g_{\mu\nu} = \frac{L^2}{z^2} \eta_{\mu\nu}, \quad g^{\mu\nu} = \frac{z^2}{L^2} \eta^{\mu\nu},
\end{equation}
其中 $\eta_{\mu\nu} = \text{diag}(-1, 1, \ldots, 1)$ 是闵可夫斯基度规。

Christoffel 符号的定义为：
\begin{equation}
    \Gamma^\mu_{\nu\rho} = \frac{1}{2} g^{\mu\sigma} (\partial_\nu g_{\rho\sigma} + \partial_\rho g_{\nu\sigma} - \partial_\sigma g_{\nu\rho}).
\end{equation}

由于度规只依赖于 $z$ 坐标，只有 $\partial_z g_{\mu\nu} \neq 0$。计算 $\partial_z g_{\mu\nu}$：
\begin{equation}
    \partial_z g_{\mu\nu} = \partial_z \left(\frac{L^2}{z^2} \eta_{\mu\nu}\right) = -\frac{2L^2}{z^3} \eta_{\mu\nu} = -\frac{2}{z} g_{\mu\nu}.
\end{equation}

现在计算各 Christoffel 符号。首先，当 $\mu = z$ 时：
\begin{equation}
    \Gamma^z_{\nu\rho} = \frac{1}{2} g^{zz} (\partial_\nu g_{\rho z} + \partial_\rho g_{\nu z} - \partial_z g_{\nu\rho}).
\end{equation}

对于 $\Gamma^z_{tt}$：
\begin{equation}
    \Gamma^z_{tt} = \frac{1}{2} g^{zz} (\partial_t g_{tz} + \partial_t g_{tz} - \partial_z g_{tt}) = \frac{1}{2} \cdot \frac{z^2}{L^2} \cdot \left(-\partial_z g_{tt}\right).
\end{equation}
由于 $g_{tt} = -L^2/z^2$，$\partial_z g_{tt} = 2L^2/z^3$，因此：
\begin{equation}
    \Gamma^z_{tt} = \frac{1}{2} \cdot \frac{z^2}{L^2} \cdot \left(-\frac{2L^2}{z^3}\right) = -\frac{1}{z}.
\end{equation}

对于 $\Gamma^z_{xx}$（$x$ 是任意边界空间坐标）：
\begin{equation}
    \Gamma^z_{xx} = \frac{1}{2} g^{zz} (-\partial_z g_{xx}) = \frac{1}{2} \cdot \frac{z^2}{L^2} \cdot \left(-\frac{2L^2}{z^3}\right) = -\frac{1}{z}.
\end{equation}

对于 $\Gamma^z_{zz}$：
\begin{equation}
    \Gamma^z_{zz} = \frac{1}{2} g^{zz} (\partial_z g_{zz} + \partial_z g_{zz} - \partial_z g_{zz}) = \frac{1}{2} g^{zz} \partial_z g_{zz} = \frac{1}{2} \cdot \frac{z^2}{L^2} \cdot \left(-\frac{2L^2}{z^3}\right) = -\frac{1}{z}.
\end{equation}

对于 $\Gamma^t_{tz}$：
\begin{equation}
    \Gamma^t_{tz} = \frac{1}{2} g^{tt} (\partial_t g_{zt} + \partial_z g_{tt} - \partial_t g_{tz}) = \frac{1}{2} g^{tt} \partial_z g_{tt} = \frac{1}{2} \cdot \left(-\frac{z^2}{L^2}\right) \cdot \frac{2L^2}{z^3} = -\frac{1}{z}.
\end{equation}

类似地，$\Gamma^x_{xz} = -1/z$。

\begin{remark}
所有非零的 Christoffel 符号都等于 $-1/z$，这反映了 AdS 空间的高度对称性。物理上，$-1/z$ 项对应于"引力势"的梯度：当我们向体空间内部运动（$z$ 增加）时，引力势变强，这导致测地线向体空间弯曲。
\end{remark}

\begin{example}[纯径向测地线]
考虑纯径向运动：$dt/d\tau = E$，$d\vec{x}/d\tau = 0$。测地线方程给出：
\begin{equation}
    \frac{d^2 z}{d\tau^2} - \frac{1}{z} \left(\frac{dz}{d\tau}\right)^2 + \frac{E^2}{z} = 0.
\end{equation}
令 $u = \ln z$，方程变为 $u'' + E^2 e^{-2u} = 0$，其解为：
\begin{equation}
    z(\tau) = z_0 e^{\pm E \tau / L}.
\end{equation}
正号对应于从边界向体空间运动（$z$ 增加），负号对应于从体空间向边界运动（$z$ 减小）。测地线在有限仿射参数内到达边界（$z \to 0$），但需要无穷大的仿射参数才能到达庞加莱视界（$z \to \infty$）。
\end{example}

\begin{example}[类空测地线的推导]
考虑类空测地线在 $t = \text{const}$ 超曲面上的投影。对于连接边界上两点 $(\vec{x}_1, z=0)$ 和 $(\vec{x}_2, z=0)$ 的测地线，我们来详细推导其形状。

\textbf{物理动机}：在 AdS 空间中，测地线是"最短路径"。由于 AdS 的负曲率，测地线会向体空间内部弯曲，形成一个"凹"的形状。这类似于在双曲面上，两点之间的最短路径是向内弯曲的弧线。

\textbf{步骤 1：建立测地线方程}。在 $t = \text{const}$ 超曲面上，度规为：
\begin{equation}
    ds^2 = \frac{L^2}{z^2} (dx^2 + dz^2).
\end{equation}
测地线方程为：
\begin{equation}
    \frac{d^2 z}{dx^2} - \frac{1}{z} \left(\frac{dz}{dx}\right)^2 + \frac{1}{z} = 0.
\end{equation}

\textbf{步骤 2：求解方程}。令 $z' = dz/dx$，方程变为：
\begin{equation}
    z'' - \frac{1}{z} (z')^2 + \frac{1}{z} = 0.
\end{equation}
由于方程不显含 $x$，存在守恒量：
\begin{equation}
    \mathcal{H} = z' \frac{\partial f}{\partial z'} - f = \text{const},
\end{equation}
其中 $f = \sqrt{1 + z'^2}/z$。计算得：
\begin{equation}
    \mathcal{H} = -\frac{1}{z\sqrt{1 + z'^2}} = -\frac{1}{R},
\end{equation}
其中 $R$ 是常数（半圆的半径）。

\textbf{步骤 3：求解形状}。从守恒量得到：
\begin{equation}
    1 + z'^2 = \frac{R^2}{z^2},
\end{equation}
即：
\begin{equation}
    z' = \pm \frac{\sqrt{R^2 - z^2}}{z}.
\end{equation}
积分得到：
\begin{equation}
    z(x) = \sqrt{R^2 - (x - x_0)^2},
\end{equation}
其中 $x_0$ 是积分常数。这是一个半圆，其最大深度为 $z_* = R$。

\textbf{物理意义}：类空测地线是半圆，这反映了 AdS 空间的负曲率。在闵可夫斯基时空中，类空测地线是直线；在 AdS 空间中，类空测地线是半圆。这个结果在全息纠缠熵的计算中起着核心作用。
\end{example}

\begin{figure}[htbp]
\centering
\begin{tikzpicture}[scale=1.2]
    % 绘制 AdS 边界
    \draw[thick, blue] (-4,0) -- (4,0) node[right] {边界};
    
    % 绘制类空测地线（半圆）
    \draw[thick, red, domain=-2:2, samples=100] plot (\x, {sqrt(4 - \x*\x)});
    
    % 标注端点
    \fill[red] (-2,0) circle (0.1) node[below] {$\vec{x}_1$};
    \fill[red] (2,0) circle (0.1) node[below] {$\vec{x}_2$};
    
    % 标注最大深度
    \draw[dashed, gray] (0,2) -- (0,0);
    \node at (0.5,2) {$z_* = R$};
    
    % 标注距离
    \draw[<->] (-2,-0.5) -- (2,-0.5);
    \node at (0,-0.8) {$L = 2R$};
    
    % 绘制径向方向
    \draw[->, gray] (3,0) -- (3,-2);
    \node at (3.3,-1) {$z$};
\end{tikzpicture}
\caption{AdS 空间中的类空测地线。连接边界上两点的测地线是半圆，其最大深度为 $z_* = R = L/2$。}
\label{fig:geodesic}
\end{figure}

\begin{figure}[htbp]
\centering
\begin{tikzpicture}[scale=1.5]
    % 绘制 Poincaré 圆盘
    \draw[thick, blue] (0,0) circle (1);
    
    % 填充
    \fill[blue!5] (0,0) circle (1);
    
    % 绘制几条测地线（半圆）
    \draw[thick, red, domain=0.1:0.9, samples=50] plot (\x, {sqrt(0.25 - (\x-0.5)*(\x-0.5))});
    \draw[thick, red, domain=-0.6:0.9, samples=50] plot (\x, {sqrt(1.2 - (\x-0.3)*(\x-0.3))});
    \draw[thick, red, domain=-0.9:-0.1, samples=50] plot (\x, {sqrt(1.2 - (\x+0.3)*(\x+0.3))});
    
    % 标注边界
    \node[blue] at (0,1.2) {AdS 边界（共形边界）};
    
    % 标注中心
    \fill[black] (0,0) circle (0.03) node[below right] {$i^0$};
    
    % 标注测地线
    \node[red] at (0.8,0.8) {测地线};
    
    % 绘制径向方向
    \draw[->, gray, thick] (0,0) -- (0.7,0.7);
    \node[gray] at (0.9,0.3) {$z$（径向）};
\end{tikzpicture}
\caption{Poincaré 圆盘模型（AdS$_2$）。蓝色圆是共形边界，红色曲线是测地线。在 AdS 空间中，测地线是与边界正交的圆弧 \cite{Witten1998}。}
\label{fig:Poincare_disk}
\end{figure}

\subsection{测地线距离与两点函数}

在 AdS/CFT 对偶中，边界 CFT 的两点关联函数可以通过体空间中的传播子来计算 \cite{Witten1998}。对于标量场，传播子正比于 $e^{-\Delta \sigma/L}$，其中 $\Delta$ 是算符的共形维数，$\sigma$ 是测地线距离。在两点距离 $|x_{12}|$ 远大于 $z$ 的极限下：
\begin{equation}
    \langle \mathcal{O}(x_1) \mathcal{O}(x_2) \rangle \sim \frac{1}{|x_{12}|^{2\Delta}},
\end{equation}
这与 CFT 中由共形对称性固定的形式一致 \cite{Maldacena1997}。

%=============================================================================
\section{AdS 中的波动方程}

在 AdS 空间中，场的行为由 Klein-Gordon 方程（对于标量场）或更一般的波动方程描述。理解这些方程的解对于计算 AdS/CFT 中的关联函数至关重要 \cite{Witten1998}。

\subsection{Klein-Gordon 方程的推导}

在 AdS 空间中，场的行为由 Klein-Gordon 方程（对于标量场）或更一般的波动方程描述。理解这些方程的解对于计算 AdS/CFT 中的关联函数至关重要 \cite{Aharony1999}。

\textbf{物理动机}：在 AdS/CFT 对偶中，边界 CFT 的算符对应于体空间中的场。为了计算关联函数，我们需要求解体空间中的波动方程。Klein-Gordon 方程是最简单的波动方程，它描述了无自旋的标量场在弯曲时空中的传播。

\textbf{从作用量出发}：标量场 $\phi$ 的作用量为：
\begin{equation}
    S = -\frac{1}{2} \int d^{d+1}x \sqrt{-g} \left( g^{\mu\nu} \partial_\mu \phi \partial_\nu \phi + m^2 \phi^2 \right).
\end{equation}
第一项是动能项，第二项是质量项。对 $\phi$ 变分，得到运动方程：
\begin{equation}
    \frac{1}{\sqrt{-g}} \partial_\mu \left( \sqrt{-g} g^{\mu\nu} \partial_\nu \phi \right) - m^2 \phi = 0.
\end{equation}
这就是弯曲时空中的 Klein-Gordon 方程：
\begin{equation}
    (\Box - m^2) \phi = 0,
\end{equation}
其中 $\Box = \frac{1}{\sqrt{-g}} \partial_\mu (\sqrt{-g} g^{\mu\nu} \partial_\nu)$ 是达朗贝尔算符。

\textbf{在庞加莱坐标下}：AdS 度规为 $ds^2 = (L^2/z^2)(-dt^2 + d\vec{x}^2 + dz^2)$，因此：
\begin{equation}
    g_{\mu\nu} = \frac{L^2}{z^2} \eta_{\mu\nu}, \quad g^{\mu\nu} = \frac{z^2}{L^2} \eta^{\mu\nu}, \quad \sqrt{-g} = \left(\frac{L}{z}\right)^{d+1}.
\end{equation}
代入 Klein-Gordon 方程：
\begin{equation}
    \frac{z^{d+1}}{L^{d+1}} \partial_\mu \left( \frac{L^{d+1}}{z^{d+1}} \frac{z^2}{L^2} \eta^{\mu\nu} \partial_\nu \phi \right) - m^2 \phi = 0.
\end{equation}
简化后得到：
\begin{equation}
    \frac{z^{d+1}}{L^2} \partial_z \left(\frac{L^2}{z^{d+1}} \partial_z \phi\right) + \frac{z^2}{L^2} \eta^{\mu\nu} \partial_\mu \partial_\nu \phi - m^2 \phi = 0.
\end{equation}

\subsection{分离变量与径向方程}

\textbf{物理动机}：由于边界是平移不变的，我们可以将场分解为不同动量的模式。这类似于量子力学中的分离变量：将多维问题分解为多个一维问题。

对方程做傅里叶变换：
\begin{equation}
    \phi(z, x) = \int \frac{d^d k}{(2\pi)^d} e^{ik \cdot x} \phi_k(z),
\end{equation}
其中 $k^2 = -\omega^2 + \vec{k}^2$。代入方程，得到径向部分满足：
\begin{equation}
    \left[\frac{1}{z^{d+1}} \partial_z \left(z^{d+1} \partial_z\right) + \frac{k^2}{z^2} - \frac{m^2 L^2}{z^2}\right] \phi_k = 0.
    \label{eq:AdS_radial}
\end{equation}
这个方程描述了场在径向方向上的行为，是 AdS/CFT 计算的核心方程。

\subsection{边界行为与质量-维数关系的推导}

\textbf{物理动机}：在边界 $z \to 0$ 附近，场的行为决定了边界 CFT 中对应算符的性质。我们需要确定场在边界附近的渐近行为，以及这如何与算符的共形维数相关联。

在边界 $z \to 0$ 附近，方程 \eqref{eq:AdS_radial} 的主导项为：
\begin{equation}
    \frac{1}{z^{d+1}} \partial_z \left(z^{d+1} \partial_z \phi_k\right) - \frac{m^2 L^2}{z^2} \phi_k \approx 0.
\end{equation}
这是一个 Euler 方程，其解具有幂律形式 $\phi_k \sim z^\alpha$。代入方程：
\begin{equation}
    \frac{1}{z^{d+1}} \partial_z \left(z^{d+1} \alpha z^{\alpha-1}\right) - \frac{m^2 L^2}{z^2} z^\alpha = 0.
\end{equation}
计算第一项：
\begin{equation}
    \frac{1}{z^{d+1}} \partial_z \left(\alpha z^{d+\alpha}\right) = \frac{1}{z^{d+1}} \cdot \alpha(d+\alpha) z^{d+\alpha-1} = \frac{\alpha(d+\alpha)}{z^2} z^\alpha.
\end{equation}
因此方程变为：
\begin{equation}
    \frac{\alpha(d+\alpha)}{z^2} z^\alpha - \frac{m^2 L^2}{z^2} z^\alpha = 0.
\end{equation}
约去 $z^\alpha/z^2$，得到特征方程：
\begin{equation}
    \alpha(d+\alpha) = m^2 L^2.
\end{equation}
定义 $\Delta = d + \alpha$（即 $\alpha = \Delta - d$），则方程变为：
\begin{equation}
    (\Delta - d) \Delta = m^2 L^2,
\end{equation}
即：
\begin{equation}
    \Delta(\Delta - d) = m^2 L^2.
    \label{eq:mass_dimension}
\end{equation}
这就是\textbf{质量-维数关系}（mass-dimension relation）！

\textbf{求解}：方程 \eqref{eq:mass_dimension} 是 $\Delta$ 的二次方程：
\begin{equation}
    \Delta^2 - d\Delta - m^2 L^2 = 0.
\end{equation}
其解为：
\begin{equation}
    \Delta_\pm = \frac{d}{2} \pm \sqrt{\frac{d^2}{4} + m^2 L^2}.
\end{equation}
两个解分别对应于场在边界附近的两种行为：
\begin{itemize}
    \item $\Delta_+$（大根）：$\phi \sim z^{\Delta_+ - d} = z^{\sqrt{d^2/4 + m^2 L^2} - d/2}$，当 $z \to 0$ 时衰减（如果 $\Delta_+ > d$）
    \item $\Delta_-$（小根）：$\phi \sim z^{\Delta_- - d} = z^{-\sqrt{d^2/4 + m^2 L^2} - d/2}$，当 $z \to 0$ 时发散
\end{itemize}

因此，场在边界附近的渐近行为为：
\begin{equation}
    \phi_k(z) \sim A(k) z^{d-\Delta} + B(k) z^\Delta,
\end{equation}
其中 $\Delta = \Delta_+$，$A(k)$ 对应于源，$B(k)$ 对应于期望值。

\begin{theorem}[质量-维数关系]
在 AdS/CFT 对偶中，体空间中质量为 $m$ 的标量场对应于边界 CFT 中共形维数为 $\Delta$ 的算符，两者满足：
\begin{equation}
    \Delta(\Delta - d) = m^2 L^2.
\end{equation}
其解为：
\begin{equation}
    \Delta_\pm = \frac{d}{2} \pm \sqrt{\frac{d^2}{4} + m^2 L^2}.
\end{equation}
通常选择 $\Delta = \Delta_+$ 作为物理的共形维数。
\end{theorem}

\begin{note}[物理意义]
质量-维数关系是 AdS/CFT 对偶中最重要的关系之一。它的物理意义是：
\begin{itemize}
    \item \textbf{质量越大 $\to$ 衰减越快}：质量越大的粒子，在体空间中传播时衰减越快，对应于边界 CFT 中维数越大的算符
    \item \textbf{无质量场}：当 $m^2 L^2 = 0$ 时，$\Delta = d$ 或 $\Delta = 0$，对应于无质量场（如规范场、应力张量）
    \item \textbf{BF 界}：当 $m^2 L^2 < 0$ 时，$\Delta$ 可能是复数，但只要满足 BF 界，$\Delta$ 仍然是实数
\end{itemize}
\end{note}

\subsection{正常izable 与 non-normalizable 模}

边界行为中的两个解具有不同的物理意义：
\begin{itemize}
    \item \textbf{Non-normalizable 模}（$z^{d-\Delta}$）：在边界上发散，对应于 CFT 中的\textbf{源}（source）。$A(k)$ 是源函数 $\phi_0(k)$。
    \item \textbf{Normalizable 模}（$z^\Delta$）：在边界上衰减，对应于 CFT 中算符的\textbf{期望值}。$B(k)$ 与 $\langle \mathcal{O}(k) \rangle$ 相关。
\end{itemize}

在 GKPW 关系中，边界条件为 $\phi \to z^{d-\Delta} \phi_0$（当 $z \to 0$），其中 $\phi_0$ 是 CFT 中的源。

\begin{remark}[GKPW 关联函数 \cite{Witten1998}]
Gubser-Klebanov-Polyakov \cite{Gubser1998} 和 Witten \cite{Witten1998} 独立地提出了 AdS/CFT 对偶中计算关联函数的基本公式（称为 GKPW 关系）。对于体空间中的标量场 $\phi$，其边界条件为 $\phi(z \to 0, x) \to z^{d-\Delta} \phi_0(x)$，其中 $\phi_0(x)$ 是边界 CFT 中算符 $\mathcal{O}(x)$ 的源。生成泛函为：
\begin{equation}
    \left\langle \exp\left(\int d^d x \, \phi_0(x) \mathcal{O}(x)\right) \right\rangle_{\rm CFT} = Z_{\rm string}[\phi \to z^{d-\Delta} \phi_0 \text{ at boundary}].
\end{equation}
在经典引力极限（$G_N \to 0$）下，弦理论的配分函数简化为 $Z_{\rm string} \approx e^{-S_{\rm on-shell}}$，因此：
\begin{equation}
    \left\langle \exp\left(\int d^d x \, \phi_0(x) \mathcal{O}(x)\right) \right\rangle_{\rm CFT} \approx e^{-S_{\rm classical}[\phi \to z^{d-\Delta} \phi_0]}.
\end{equation}
对 $\phi_0$ 求泛函导数，可以得到 CFT 中算符的关联函数。例如，两点函数为：
\begin{equation}
    \langle \mathcal{O}(x_1) \mathcal{O}(x_2) \rangle = \frac{(2\Delta - d) \Gamma(\Delta)}{\pi^{d/2} \Gamma(\Delta - d/2)} \cdot \frac{1}{|x_{12}|^{2\Delta}},
\end{equation}
这与 CFT 中由共形对称性固定的形式一致 \cite{Witten1998}。
\end{remark}

\subsection{Breitenlohner-Freedman 界}

对于 AdS 空间中的标量场，存在一个质量下界，称为\textbf{Breitenlohner-Freedman（BF）界}：
\begin{equation}
    m^2 L^2 \geq -\frac{d^2}{4}.
\end{equation}
当 $m^2 L^2 < 0$ 时，场是\textbf{快子}（tachyonic）的，但只要满足 BF 界，AdS 的负曲率可以稳定这个不稳定性。

\begin{remark}
BF 界的物理意义是：AdS 空间的负曲率提供了一个"引力势阱"，可以束缚快子场。这与闵可夫斯基时空不同——在闵可夫斯基时空中，任何负质量平方的场都是不稳定的。BF 界在 AdS/CFT 对偶中有重要的物理含义：它对应于边界 CFT 中算符维数的么正性界 \cite{Aharony1999}。
\end{remark}

\begin{proof}[BF 界的推导 \cite{Aharony1999}]
我们来证明：当 $m^2 L^2 \geq -d^2/4$ 时，标量场在 AdS 空间中是稳定的。考虑标量场的作用量：
\begin{equation}
    S = -\frac{1}{2} \int d^{d+1}x \sqrt{-g} \left( g^{\mu\nu} \partial_\mu \phi \partial_\nu \phi + m^2 \phi^2 \right).
\end{equation}
在全局坐标下，$\sqrt{-g} = (L/\cos\rho)^{d+1}$。对作用量做分部积分：
\begin{equation}
    S = -\frac{1}{2} \int d^{d+1}x \sqrt{-g} \left( -\phi \Box \phi + m^2 \phi^2 \right) + \text{边界项}.
\end{equation}
对于模式 $\phi \sim e^{-i\omega t} Y_{\ell}(\Omega) f(\rho)$，能量为 $E = \omega$。如果所有模式都有 $E^2 \geq 0$（即 $\omega$ 为实数），则系统是稳定的。

在 AdS 边界附近（$\rho \to \pi/2$），径向方程的渐近行为为 $f \sim \cos^\alpha\rho$，其中 $\alpha(\alpha-d) = m^2 L^2$。为了使作用量有下界（即能量有下界），我们需要径向方程没有负范数态。这要求：
\begin{equation}
    \alpha = \frac{d}{2} + \sqrt{\frac{d^2}{4} + m^2 L^2} \in \mathbb{R},
\end{equation}
即 $m^2 L^2 \geq -d^2/4$。当 $m^2 L^2 = -d^2/4$ 时，$\alpha = d/2$，这是临界情况。因此，BF 界为：
\begin{equation}
    m^2 L^2 \geq -\frac{d^2}{4}.
\end{equation}
\end{proof}

\subsection{精确解：AdS$_{5}$ 中的标量场}

作为具体例子，考虑 AdS$_5$（$d=4$）中的标量场。径向方程 \eqref{eq:AdS_radial} 变为：
\begin{equation}
    \left[\frac{1}{z^5} \partial_z \left(z^5 \partial_z\right) - k^2 - \frac{m^2 L^2}{z^2}\right] \phi_k = 0.
\end{equation}
令 $\phi_k(z) = z^2 f(z)$，方程变为修正的贝塞尔方程：
\begin{equation}
    z^2 f'' + z f' - \left(\Delta^2 + k^2 z^2\right) f = 0,
\end{equation}
其中 $\Delta = 2 + \sqrt{4 + m^2 L^2}$。解为：
\begin{equation}
    \phi_k(z) = z^2 \left[ A(k) K_\Delta(|k|z) + B(k) I_\Delta(|k|z) \right],
\end{equation}
其中 $K_\Delta$ 和 $I_\Delta$ 是修正的贝塞尔函数。在边界 $z \to 0$ 附近：
\begin{equation}
    \phi_k(z) \sim A(k) \Gamma(\Delta) \left(\frac{|k|z}{2}\right)^{-\Delta} 2^{\Delta-1} + B(k) \Gamma(-\Delta) \left(\frac{|k|z}{2}\right)^\Delta 2^{-\Delta-1}.
\end{equation}
这与一般分析一致：$A(k)$ 对应于源，$B(k)$ 对应于期望值。

%=============================================================================
\section{AdS 的热力学：黑洞解}

AdS 空间可以支持黑洞解。与渐近平坦时空中的黑洞不同，AdS 黑洞具有独特的热力学性质，这在 AdS/CFT 对偶中对应于边界 CFT 的热态 \cite{Hawking1983, Maldacena1997}。

\subsection{AdS-Schwarzschild 黑洞}

\begin{definition}[AdS-Schwarzschild 黑洞]
AdS$_{d+1}$ 中的 Schwarzschild 黑洞度规为：
\begin{equation}
    ds^2 = -f(r) dt^2 + \frac{dr^2}{f(r)} + r^2 d\Omega_{d-1}^2,
    \label{eq:AdS_Schwarzschild}
\end{equation}
其中：
\begin{equation}
    f(r) = 1 + \frac{r^2}{L^2} - \frac{\mu}{r^{d-2}}.
\end{equation}
参数 $\mu$ 与黑洞质量 $M$ 的关系为：
\begin{equation}
    M = \frac{(d-1) \omega_{d-1}}{16\pi G_N} \mu,
\end{equation}
其中 $\omega_{d-1} = 2\pi^{d/2}/\Gamma(d/2)$ 是单位 $(d-1)$ 维球面的面积。
\end{definition}

\begin{note}[与渐近平坦黑洞的区别]
AdS 黑洞与渐近平坦时空中的 Schwarzschild 黑洞有本质区别：
\begin{itemize}
    \item \textbf{无毛定理}：AdS 黑洞同样满足无毛定理，但其渐近行为不同
    \item \textbf{温度行为}：AdS 黑洞的温度有最小值，而渐近平坦黑洞的温度随质量减小而增大
    \item \textbf{稳定性}：AdS 黑洞是热力学稳定的，而小的渐近平坦黑洞会通过霍金辐射蒸发
\end{itemize}
\end{note}

\subsection{视界与霍金温度}

视界半径 $r_h$ 由 $f(r_h) = 0$ 确定：
\begin{equation}
    1 + \frac{r_h^2}{L^2} = \frac{\mu}{r_h^{d-2}}.
\end{equation}
对于给定的质量 $M$，这个方程可能有一个或两个解。大视界 $r_+$ 对应于黑洞，小视界 $r_-$ 对应于 Cauchy 视界。

\textbf{霍金温度}（Hawking temperature）通过表面引力 $\kappa$ 计算：
\begin{equation}
    T = \frac{\kappa}{2\pi} = \frac{f'(r_h)}{4\pi} = \frac{1}{4\pi} \left( \frac{d \, r_h}{L^2} + \frac{d-2}{r_h} \right).
    \label{eq:Hawking_T}
\end{equation}

\begin{figure}[htbp]
\centering
\begin{tikzpicture}[scale=1.2]
    % 绘制坐标轴
    \draw[thick, ->] (0,0) -- (5,0) node[right] {$r_h$};
    \draw[thick, ->] (0,0) -- (0,4) node[above] {$T$};
    
    % 绘制温度曲线
    \draw[thick, blue, domain=0.5:4.5, samples=100] 
        plot (\x, {0.5*\x + 1/\x});
    
    % 标注最小值
    \fill[red] (1.414, 1.414) circle (0.1) node[above right] {$T_{\min}$};
    
    % 标注渐近行为
    \draw[dashed, gray] (0,0) -- (4,2) node[midway, above left] {$T \sim r_h$};
    \draw[dashed, gray] (0.5,2) -- (0,4) node[midway, left] {$T \sim 1/r_h$};
    
    % 标注 AdS 半径
    \draw[dashed, red] (2,0) -- (2,2.5) node[above] {$r_h \sim L$};
\end{tikzpicture}
\caption{AdS 黑洞的温度 $T$ 与视界半径 $r_h$ 的关系。温度在 $r_h \sim L$ 处达到最小值 $T_{\min} \sim d/(2\pi L)$。}
\label{fig:AdS_BH_T}
\end{figure}

\begin{property}[AdS 黑洞的温度行为]
与渐近平坦时空中的 Schwarzschild 黑洞不同，AdS 黑洞的温度有最小值：
\begin{itemize}
    \item 当 $r_h \gg L$ 时，$T \approx d r_h/(4\pi L^2)$，温度与视界半径成正比
    \item 当 $r_h \ll L$ 时，$T \approx (d-2)/(4\pi r_h)$，温度与视界半径成反比
    \item 温度在 $r_h \sim L$ 处达到最小值 $T_{\min} \sim d/(2\pi L)$
\end{itemize}
\end{property}

\subsection{黑洞热力学}

AdS 黑洞具有丰富的热力学性质。我们来详细讨论这些性质。

\textbf{贝肯斯坦-霍金熵}：黑洞的熵由贝肯斯坦-霍金公式给出：
\begin{equation}
    S = \frac{A}{4G_N} = \frac{\omega_{d-1} r_h^{d-1}}{4G_N}.
\end{equation}
这个熵是广延量，与边界 CFT 的热力学熵一致。

\textbf{热力学第一定律}：AdS 黑洞的热力学量满足热力学第一定律：
\begin{equation}
    dM = T dS.
\end{equation}
将 $M$、$T$ 和 $S$ 表示为 $r_h$ 的函数，可以验证这个关系：
\begin{align}
    dM &= \frac{(d-1) \omega_{d-1}}{16\pi G_N} \left( (d-2) r_h^{d-3} + d \frac{r_h^{d-1}}{L^2} \right) dr_h, \\
    T dS &= \frac{1}{4\pi} \left( \frac{d \, r_h}{L^2} + \frac{d-2}{r_h} \right) \cdot \frac{(d-1) \omega_{d-1} r_h^{d-2}}{4G_N} dr_h,
\end{align}
两者相等。

\textbf{自由能}：黑洞的亥姆霍兹自由能为：
\begin{equation}
    F = M - TS = -\frac{\omega_{d-1} r_h^d}{16\pi G_N L^2}.
\end{equation}
在高温极限 $r_h \gg L$ 下：
\begin{equation}
    F \approx -\frac{\pi^2}{2} N^2 T^4 V,
\end{equation}
其中 $V$ 是边界体积。这与 $\mathcal{N}=4$ SYM 在强耦合极限下的自由能一致。

\textbf{霍金-佩奇相变}：在 AdS 空间中，热 AdS（thermal AdS）和 AdS 黑洞之间发生相变，称为\textbf{霍金-佩奇相变}（Hawking-Page transition）\cite{Hawking1983}。当温度 $T > T_c$ 时，AdS 黑洞是稳定的状态；当 $T < T_c$ 时，热 AdS 是稳定的状态。临界温度为：
\begin{equation}
    T_c = \frac{d}{2\pi L}.
\end{equation}
在临界温度处，自由能连续，但其一阶导数（熵）不连续，因此这是一阶相变。

\begin{remark}
霍金-佩奇相变在 AdS/CFT 对偶中对应于边界 CFT 中的\textbf{禁闭/解禁闭相变}（confinement/deconfinement transition）\cite{Witten1998b}。在禁闭相（$T < T_c$）中，自由能与 $N^0$ 成正比，对应于强子相；在解禁闭相（$T > T_c$）中，自由能与 $N^2$ 成正比，对应于夸克-胶子等离子体相 \cite{Aharony1999}。这个相变是理解 QCD 相图的重要工具。
\end{remark}

\begin{remark}[AdS 黑洞与全息对偶 \cite{Aharony1999}]
AdS 黑洞在全息对偶中有深刻的物理意义 \cite{Aharony1999}：
\begin{itemize}
    \item \textbf{热态对应}：AdS 黑洞对应于边界 CFT 的热态。黑洞的温度等于 CFT 的温度，黑洞的熵等于 CFT 的热力学熵
    \item \textbf{大 $N$ 极限}：在 $N \to \infty$ 极限下，黑洞的量子涨落可以忽略，热力学量与 $N^2$ 成正比，这与 $\mathcal{N}=4$ SYM 的自由度数一致
    \item \textbf{剪切粘滞系数}：AdS 黑洞的准正则模（quasinormal modes）对应于边界 CFT 中的流体力学模式，剪切粘滞系数与熵密度之比为 $\eta/s = 1/(4\pi)$，这是一个普适的结果 \cite{Policastro2001, Kovtun2005}
\end{itemize}
\end{remark}

%=============================================================================
\section{AdS 空间的直观理解}

为了帮助读者建立对 AdS 空间的直观理解，我们总结几个关键的物理图像。

\subsection{几何直观}

\textbf{"碗"状几何}：AdS 空间可以被想象为一个"碗"状的几何。边界在碗的边缘（$z = 0$），而碗的底部（$z = \infty$）是 AdS 的"中心"。径向坐标 $z$ 对应于从边界到体空间的"深度"。

这个"碗"的比喻有一个重要的特点：碗的"体积"在边界附近发散。这是因为度规因子 $L^2/z^2$ 在 $z \to 0$ 时发散。这意味着 AdS 空间在边界附近有无限大的"体积"，这与 CFT 中 UV 发散的物理图像一致。

\textbf{径向方向作为能量尺度}：在 AdS/CFT 对偶中，径向坐标 $z$ 可以被理解为能量尺度的倒数：
\begin{equation}
    z \sim \frac{1}{\mu},
\end{equation}
其中 $\mu$ 是能量尺度。边界 $z \to 0$ 对应于紫外（UV，高能），体空间内部 $z \to \infty$ 对应于红外（IR，低能）。这个对应关系是 AdS/CFT 中重整化群流的几何实现。

\textbf{AdS 与双曲空间}：AdS 空间与双曲空间（hyperbolic space）有密切联系。AdS$_{d+1}$ 的空间部分是 $d$ 维双曲空间 $H^d$。双曲空间是常负曲率的黎曼流形，其度规为：
\begin{equation}
    ds^2_{H^d} = \frac{L^2}{z^2} \left( d\vec{x}^2 + dz^2 \right).
\end{equation}
这恰好是 AdS 度规的空间部分。双曲空间的等距群是 $SO(1,d)$，这对应于 AdS 空间等距群 $SO(2,d)$ 的子群。

%=============================================================================
\section{AdS 空间的物理意义与应用}

AdS 空间不仅仅是一个数学构造，它在理论物理学中有着深刻的物理意义和广泛的应用。我们来详细讨论这些方面 \cite{Maldacena1997}。

\subsection{AdS 空间与量子引力}

AdS 空间在量子引力研究中扮演着核心角色。这是因为：
\begin{itemize}
    \item \textbf{负曲率的稳定性}：AdS 空间的负曲率提供了一个"引力势阱"，可以束缚场的涨落。这使得 AdS 空间在量子引力中是稳定的，而渐近平坦时空可能不是。
    \item \textbf{边界条件}：AdS 空间具有类时的边界，这为定义量子引力理论提供了自然的边界条件。我们可以指定边界上的场值，这对应于在 CFT 中添加源。
    \item \textbf{全息性质}：AdS 空间的全息性质（即体空间的物理可以由边界理论描述）为量子引力提供了一个非微扰的定义 \cite{Maldacena1997}。
\end{itemize}

\begin{remark}[AdS/CFT 字典 \cite{Maldacena1997}]
AdS/CFT 对偶建立了一套精确的对应关系（称为"字典"）\cite{Maldacena1997}：
\begin{itemize}
    \item \textbf{态的对应}：CFT 的真空态对应于纯 AdS 空间；CFT 的热态对应于 AdS 黑洞 \cite{Hawking1983}
    \item \textbf{算符的对应}：CFT 中的算符 $\mathcal{O}$ 对应于体空间中的场 $\phi$，两者通过质量-维数关系 $m^2 L^2 = \Delta(\Delta-d)$ 相联系
    \item \textbf{对称性的对应}：CFT 的共形群 $SO(2,d)$ 对应于 AdS 的等距群 $SO(2,d)$
    \item \textbf{耦合常数的对应}：CFT 的 't Hooft 耦合 $\lambda = g_{\rm YM}^2 N$ 对应于弦理论参数 $L^4/\alpha'^2$
    \item \textbf{自由度的对应}：CFT 的中心荷 $c \sim N^2$ 对应于引力理论中 $L^{d-1}/(8\pi G_N)$
\end{itemize}
这个字典使得我们可以通过计算体空间中的经典引力问题来得到强耦合 CFT 的非微扰信息。
\end{remark}

\begin{figure}[htbp]
\centering
\begin{tikzpicture}[scale=1.0]
    % 绘制 AdS 体空间（碗状）
    \draw[thick, blue, domain=-3:3, samples=100] plot (\x, {-0.1*\x*\x + 3});
    \draw[thick, blue] (-3,2.1) -- (-3,4);
    \draw[thick, blue] (3,2.1) -- (3,4);
    \draw[thick, blue] (-3,4) -- (3,4);
    
    % 填充
    \fill[blue!10] (-3,2.1) .. controls (-1.5,2.5) and (1.5,2.5) .. (3,2.1) -- (3,4) -- (-3,4) -- cycle;
    
    % 标注边界
    \draw[thick, red] (-3,4) -- (3,4) node[midway, above] {CFT 边界};
    
    % 标注径向方向
    \draw[->, thick, gray] (0,4) -- (0,2.5) node[midway, right] {$z$};
    
    % 标注体空间
    \node at (0,3.2) {AdS 体空间};
    
    % 绘制场线
    \draw[->, green!60!black, thick] (-1.5,3.8) .. controls (-1,3) and (1,3) .. (1.5,3.8);
    \draw[->, green!60!black, thick] (-2,3.8) .. controls (-1.5,2.8) and (1.5,2.8) .. (2,3.8);
    
    % 标注场
    \node[green!60!black] at (0,2.6) {$\phi$};
    
    % 标注算符
    \node[red] at (-2,4.3) {$\mathcal{O}_1$};
    \node[red] at (2,4.3) {$\mathcal{O}_2$};
    
    % 标注对应关系
    \draw[<->, thick, purple] (-2.5,3.5) -- (-2.5,4.5);
    \node[purple, left] at (-2.5,4) {$\Delta(\Delta-d)=m^2L^2$};
\end{tikzpicture}
\caption{AdS/CFT 对偶的示意图。体空间中的标量场 $\phi$（绿色）对应于边界 CFT 中的算符 $\mathcal{O}$（红色），两者通过质量-维数关系联系 \cite{Maldacena1997}。}
\label{fig:AdS_CFT_dictionary}
\end{figure}

\subsection{AdS 空间与黑洞物理}

AdS 空间中的黑洞物理与渐近平坦时空中的黑洞物理有本质区别：
\begin{itemize}
    \item \textbf{热力学稳定性}：AdS 黑洞是热力学稳定的，而小的渐近平坦黑洞会通过霍金辐射蒸发 \cite{Hawking1975}。这是因为 AdS 空间的"盒子"效应限制了辐射的逃逸。
    \item \textbf{霍金-佩奇相变}：在 AdS 空间中，热 AdS 和 AdS 黑洞之间发生相变，这对应于边界 CFT 中的禁闭/解禁闭相变 \cite{Witten1998b}。这个相变是理解 QCD 相图的重要工具。
    \item \textbf{黑洞信息悖论}：AdS/CFT 对偶为黑洞信息悖论提供了新的视角。由于边界 CFT 是幺正的，黑洞蒸发过程必须保持信息守恒 \cite{Maldacena1997}。
\end{itemize}

\subsection{AdS 空间与凝聚态物理}

AdS 空间在凝聚态物理中也有重要应用：
\begin{itemize}
    \item \textbf{全息超导体}：通过在 AdS 黑洞背景中引入复标量场，可以模拟超导相变 \cite{Hartnoll2007}。当温度低于临界温度时，标量场凝聚，破坏 $U(1)$ 对称性，系统进入超导相。
    \item \textbf{全息非费米液体}：AdS/CFT 可以描述具有非费米液体行为的强耦合系统 \cite{Liu2006}。这些系统无法用传统的朗道费米液体理论描述。
    \item \textbf{全息量子临界点}：AdS/CFT 可以描述量子相变的临界行为 \cite{Hartnoll2009}。量子临界点的临界指数可以通过求解体空间中的场方程得到。
\end{itemize}

\subsection{AdS 空间与宇宙学}

虽然 AdS 空间不是我们宇宙的好的近似（我们的宇宙具有正宇宙学常数），但它在宇宙学研究中也有应用：
\begin{itemize}
    \item \textbf{AdS/CFT 与宇宙学}：AdS/CFT 对偶的思想可以推广到宇宙学设置。例如，de Sitter/CFT 对偶试图为具有正宇宙学常数的时空建立全息描述 \cite{Maldacena1998}。
    \item \textbf{暴胀模型}：某些暴胀模型可以在 AdS 空间中实现。AdS 的负曲率可以驱动宇宙的加速膨胀。
    \item \textbf{暗能量}：AdS/CFT 对偶为理解暗能量提供了新的视角。暗能量可能与真空能量有关，而 AdS/CFT 对偶可以帮助我们理解真空能量的性质。
\end{itemize}

\begin{remark}[大 $N$ 极限与 $1/N$ 展开 \cite{Aharony1999}]
AdS/CFT 对偶的一个重要方面是它在大 $N$ 极限下的行为 \cite{Aharony1999}。对于 $SU(N)$ 规范理论：
\begin{itemize}
    \item \textbf{领头阶}（$N \to \infty$，$\lambda$ 固定）：规范理论的关联函数由平面图（planar diagrams）主导，对应于体空间中的经典引力
    \item \textbf{弦修正}（$1/\alpha'$ 展开）：对应于规范理论中的 $\alpha'$ 修正，即高导数算符的贡献
    \item \textbf{量子修正}（$G_N$ 展开，即 $1/N^2$ 展开）：对应于体空间中的弦圈图，即规范理论中的非平面图贡献
\end{itemize}
$1/N$ 展开的拓扑展开性质使得 AdS/CFT 对偶具有深刻的数学结构：规范理论的微扰展开按照 Feynman 图的亏格（genus）分类，这与弦理论的弦耦合常数 $g_s \sim 1/N$ 的展开一致。
\end{remark}

\subsection{AdS 空间的历史}

AdS 空间的历史可以追溯到 19 世纪：
\begin{itemize}
    \item \textbf{1883 年}：Wilhelm Killing 在研究非欧几里得几何时，首次发现了 AdS 空间。
    \item \textbf{1917 年}：爱因斯坦在研究宇宙学时，考虑了具有负宇宙学常数的时空。
    \item \textbf{1983 年}：Hawking 和 Page 发现了 AdS 黑洞的热力学性质，这为后来的 AdS/CFT 对偶奠定了基础。
    \item \textbf{1997 年}：Maldacena 提出了 AdS/CFT 对偶，将 AdS 空间与共形场论联系起来 \cite{Maldacena1997}。这是理论物理学中最深刻的发现之一。
    \item \textbf{1998 年}：Gubser-Klebanov-Polyakov \cite{Gubser1998} 和 Witten \cite{Witten1998} 独立地给出了 AdS/CFT 对偶中计算关联函数的具体公式（GKPW 关系）。
    \item \textbf{1999 年}：Aharony、Gubser、Maldacena、Ooguri 和 Oz 发表了关于 AdS/CFT 对偶的综合性综述 \cite{Aharony1999}，系统地总结了大 $N$ 规范理论与弦理论的对偶关系。
\end{itemize}

\section{小结}

在本节中，我们系统地介绍了反德西特空间的几何性质 \cite{Aharony1999}：
\begin{itemize}
    \item AdS 空间是具有常负曲率的最大对称时空，可以通过嵌入在 $\mathbb{R}^{2,d}$ 中的超曲面定义
    \item AdS 空间有多种坐标系：全局坐标覆盖整个空间，庞加莱坐标覆盖一半，Fefferman-Graham 坐标适用于全息重整化 \cite{Witten1998}
    \item AdS 空间的等距群 $SO(2,d)$ 与 $d$ 维共形群同构，这是 AdS/CFT 对偶的对称性基础 \cite{Maldacena1997}
    \item AdS 空间的因果结构独特：类光无穷远是类时的，光信号可以在有限时间内往返边界 \cite{Witten1998}
    \item 测地线在 AdS 空间中的行为为全息纠缠熵的计算提供了基础
    \item Klein-Gordon 方程的解具有边界行为 $z^{d-\Delta}$ 和 $z^\Delta$，分别对应于源和期望值 \cite{Witten1998}
    \item AdS 黑洞具有有限温度和熵，霍金-佩奇相变对应于禁闭/解禁闭相变 \cite{Hawking1983, Witten1998b}
    \item AdS 空间在量子引力、黑洞物理、凝聚态物理和宇宙学中都有重要应用
\end{itemize}
这些几何性质为理解 AdS/CFT 对偶提供了坚实的基础。在下一节中，我们将介绍共形场论的基本概念，这是 AdS/CFT 对偶中边界理论的基础 \cite{Maldacena1997}。
